The political economy of redistricting reform is characterized by a fundamental collective action problem: the legislators who must enact reform are the same officials whose electoral survival depends on avoiding it. A formal model of incumbency rent---the electoral cushion provided by favorable district lines---explains why reform through ordinary legislative channels fails at a rate of nearly 90\%, even when public opinion supports reform by wide margins.

Three bypass channels circumvent this problem. Direct democracy through ballot initiatives succeeds at 73\%, constrained primarily by the 26-state availability of the initiative mechanism. Judicial intervention succeeds when plaintiffs can demonstrate constitutional or statutory violations, but proceeds slowly. Federal legislation under Article I, \S4 is theoretically the most powerful channel, but has not yet been used for redistricting reform in the modern era.

The model generates testable predictions that hold in the data: reform probability increases with electoral competitiveness ($p = 0.042$) and with the interaction of initiative availability and competitive conditions (78\% success versus 14\%). These findings are consistent with the theory that competitiveness reduces the value of incumbency rents and that bypass channels succeed precisely because they do not require incumbent consent.

The implications for this research program are direct. The \texttt{bisect} algorithmic redistricting platform solves the technical problem: it produces neutral, compact, legally defensible maps at low cost (P.5). But technical solutions require adoption pathways. Papers P.1 through P.5 each develop one pathway in sufficient detail to be actionable: federal statutory text (P.1), ballot initiative design (P.2), commission integration protocols (P.3), court-order templates (P.4), and cost--benefit analysis (P.5). The choice among pathways is a political economy question; this paper provides the framework for making it.

The core insight is that reform does not require persuading incumbents that fairness matters. It requires choosing pathways where incumbents do not control the outcome. Voters, courts, and Congress each possess the institutional authority to mandate neutral redistricting; each faces different procedural requirements; and each has precedent supporting algorithmic solutions. The toolkit is complete. The political conditions for using it vary by state and by federal political cycle. Reform advocates should calibrate their strategy accordingly.
